\documentclass[11pt,a4paper]{article}

\usepackage[margin=21mm]{geometry}

\usepackage{graphicx,booktabs,longtable,amsmath,amssymb}

\usepackage[T1]{fontenc}

\usepackage{hyperref}

\hypersetup{colorlinks=true,linkcolor=blue,urlcolor=blue}

\setlength{\parskip}{5pt}

\setlength{\parindent}{0pt}

\begin{document}

\title{Experiment 1 on 28-10-2026\\[0.35em]\large learned post-training calibration scaling of spatial regularisation maps\\[0.5em]\normalsize Full technical report: four-image sweep, 50-image Approach A training, final CIGS evaluation}

\author{Reproducible PCM U-TV experimental report}

\date{28 October 2026 13:37 UTC}

\maketitle

\tableofcontents

\clearpage

\section{Executive summary: the two experiments and their distinct roles}

The first stage is an exploratory sweep of the frozen spatial U-TV map on four synthetic CIGS images. It is executed BEFORE training the scale head, to quantify the potential PSNR gains and the line-contrast trade-off. On its three exploratory noise seeds per image, the posthoc discrete-scale oracle gains 4.004 dB over c=1; the calibration-chosen pooled constant gains 4.004 dB. The residual oracle opportunity beyond that global constant is 0.000 dB. These figures are non-deployable posthoc upper bounds where applicable.

In Stage II we selected exactly 50 OTHER training images (10 from each of five PCM categories), plus 25 image-disjoint validation images (5 per category), without using the four CIGS samples for head learning or validation. The supervised labels were built from complete offline quality-vs-scale curves on a fixed 19-point logarithmic grid. We selected a small Approach A predictor by validation PSNR, then evaluated it on the four CIGS images at three NEW predeclared noise seeds. The selected model was near\_optimal\_set\_regression. Its mean test PSNR gain over c=1 was 3.964 dB, its gain over the best validation-chosen constant was -0.018 dB, and the remaining posthoc test oracle opportunity beyond the constant was 0.000 dB.

\section{Scientific aim and what is actually learned}

For noisy PCM measurements $\boldsymbol y=A\boldsymbol x^\star+\boldsymbol\varepsilon$, the pretrained U-TV model predicts a fixed positive spatial regularisation map $\boldsymbol\Lambda_\theta(\boldsymbol y)$. A *scalar* correction is introduced without altering the map-network parameters: \[\boldsymbol\Lambda_{\rm eff}(\boldsymbol y)=c_\phi(\boldsymbol y)\boldsymbol\Lambda_\theta(\boldsymbol y),\qquad \hat{\boldsymbol x}=\operatorname{PDHG}_K(\boldsymbol y,\boldsymbol\Lambda_{\rm eff}(\boldsymbol y)).\] Only $c_\phi$ is learned during Stage II; no gradient passes through PDHG during head training. For a fixed predicted nonnegative weight map the TV objective is convex in the reconstruction, although a finite-depth PDHG output need not be its exact minimiser.

\section{Stage I: exploratory four-image scale sweep (before training)}

The four synthetic CIGS image IDs are 0674, 0985, 0917 and 0864. Each has two exploratory calibration and three exploratory held-out noise realisations; their full images and line annotations were inspected before running the sweep. The same original measured data and predicted map are used for each scale within a datum. All Stage I scales, refinement centres and posthoc oracles are diagnostic; none of the resulting four-image labels were used for training or selecting the Stage II scale head.

\begin{center}
\small
\begin{tabular}{lr}
\toprule
Stage I, dB & Average PSNR gain vs c=1\\
\midrule
Pooled calibration constant & 4.004\\
Leave-one-image-out fixed constant & 4.004\\
Image-specific calibration & 4.003\\
Posthoc per-datum oracle & 4.004\\
Oracle - pooled constant & 0.000\\
\bottomrule
\end{tabular}
\end{center}

\begin{figure}[p]
\centering
\includegraphics[width=0.98\textwidth,height=0.78\textheight,keepaspectratio]{\detokenize{figures/stage1_quality_vs_scale.png}}
\caption{Stage I: PSNR and line contrast versus multiplier. Every curve is based on the four-image exploratory run and common measurement noise seeds per image. A scale yielding a higher global PSNR can remove narrow dark lines. The plots do not imply good transfer to unseen image structures.}
\end{figure}
\clearpage

\begin{figure}[p]
\centering
\includegraphics[width=0.98\textwidth,height=0.78\textheight,keepaspectratio]{\detokenize{figures/stage1_quality_line_tradeoff.png}}
\caption{Stage I: global image quality versus local line preservation. Every point is a scalar multiplier of a fixed learned parameter map, with calibration-only measurements used for selection. The diagram makes the PSNR-versus-contrast conflict visible.}
\end{figure}
\clearpage

\begin{figure}[p]
\centering
\includegraphics[width=0.98\textwidth,height=0.78\textheight,keepaspectratio]{\detokenize{figures/stage1_oracle_gap.png}}
\caption{Stage I: how much can per-datum oracle tuning add?. Only the extra gain of the posthoc oracle over one fixed calibration-chosen multiplier motivates training a per-datum scale predictor. An oracle determined using the unknown ground truth is not deployable.}
\end{figure}
\clearpage

\section{Stage II: input-data partition, noise realisations and targets}

We use exactly 10 OTHER images for training from each category CIGS, Organic, Perovskite, Reference and Silicon; total 50 training image identities. The pretrained U-TV mapping is frozen and its original train/validation split is respected. Additionally, 5 independent validation image identities per category (25 total) were reserved exclusively for choosing the scalar-head model and the global-constant baseline. The final evaluation uses the four CIGS identities from Stage I but NEW noise seeds 5101/5102/5103, excluded entirely from head training and validation.

\begin{center}
\small
\begin{tabular}{lrr}
\toprule
Category & Train images & Validation images\\
\midrule
CIGS & 10 & 5\\
Organic & 10 & 5\\
Perovskite & 10 & 5\\
Reference & 10 & 5\\
Silicon & 10 & 5\\
\bottomrule
\end{tabular}
\end{center}

The 50 train image identities each provide 3 noisy data: 150 supervised oracle-curve labels. The 25 validation images supply 50 separate noisy examples. Final test data count is 12. Repeated noise realisations on the same image are NOT independent image structures.

All target scales belong to a 19-point logarithmic grid, from $2^{-3}$ to $2^6$ in half-octave increments. For each noisy datum we retain the entire PSNR-versus-scale curve, the best sampled scale, the discrete near-optimal set within 0.1 dB of best, and a quality-weighted soft target with temperature 0.15 dB. These oracle labels are created with training/validation ground truth; the predictor sees only measurement, CG image-input and frozen-map summary features. Image IDs, ground-truth contrasts and true noise norms are EXCLUDED from the head input.

\section{Three supervised Approach A head losses and selection}

Each head is a small MLP applied to the same standardised input features. Let $z=\log_2 c$, and $z_i^\star$ be the offline best sampled scale. The baseline head minimises $\sum_i(z_\phi(\boldsymbol y_i)-z_i^\star)^2$. The near-optimal-set head minimises the squared distance from $z_\phi$ to the discrete set of scales within 0.1 dB of the oracle PSNR. The soft-classification head predicts weights over scales by matching $p^\star_{ij}\propto\exp((Q_{ij}-\max_\ell Q_{i\ell})/T)$ with $T=0.15$ dB. All head models are trained on CPU without differentiating through PDHG. Continuous predictions are quantised to the nearest common discrete grid scale for fair cached-solver evaluation.

The chosen model is determined by the highest VALIDATION reconstructed PSNR (not the smallest squared error on oracle log-scales). No final-test PSNR or Stage I test-oracle values were used to select it. Two constant baselines are frozen: one chosen across all five validation categories and one chosen solely on the independent CIGS validation images.

\begin{center}
\small
\begin{tabular}{lrr}
\toprule
Head / baseline & Selected epoch & Validation PSNR dB\\
\midrule
near optimal set regression & 70 & 32.316\\
hard log2 regression & 73 & 32.314\\
soft classification & 37 & 32.307\\
Global constant c=8.000 & -- & 32.119\\
CIGS-val constant c=8.000 & -- & 32.129\\
\bottomrule
\end{tabular}
\end{center}

\begin{figure}[p]
\centering
\includegraphics[width=0.98\textwidth,height=0.78\textheight,keepaspectratio]{\detokenize{figures/head_validation_psnr.png}}
\caption{Validation reconstruction PSNR during offline head training. The validation images are separate from the 50 training identities and the four CIGS final-test images. Model selection and early stopping used only these validation scores, with no PDHG backpropagation.}
\end{figure}
\clearpage

\section{Final test: PSNR gains, quality and line survival}

The four CIGS image identities are used as a small final test for the scalar head, with NEW noise seeds 5101, 5102 and 5103 for every image. Both validation constants and the chosen model were frozen first. The per-datum best scale is computed afterwards as a non-deployable test oracle. Metrics are averaged within each image across noise seeds and then equally over the four images.

\begin{center}
\small
\begin{tabular}{lrrrrr}
\toprule
Image & c=1 & All-val & CIGS-val & Learned & Oracle\\
\midrule
CIGS 0674 & 28.537 & 32.786 & 32.786 & 32.786 & 32.786\\
CIGS 0985 & 28.047 & 31.863 & 31.863 & 31.863 & 31.863\\
CIGS 0917 & 28.078 & 32.115 & 32.115 & 32.115 & 32.115\\
CIGS 0864 & 28.330 & 32.152 & 32.152 & 32.082 & 32.152\\
\bottomrule
\end{tabular}
\end{center}

\begin{center}
\small
\begin{tabular}{lr}
\toprule
Summary of four image means & PSNR gain dB\\
\midrule
Learned versus c=1 & 3.964\\
Learned versus global c & -0.018\\
Learned versus CIGS-validation c & -0.018\\
Oracle versus global c & 0.000\\
\bottomrule
\end{tabular}
\end{center}

\begin{figure}[p]
\centering
\includegraphics[width=0.98\textwidth,height=0.78\textheight,keepaspectratio]{\detokenize{figures/final_test_global_and_line.png}}
\caption{Final-test reconstruction fidelity and thin-line contrast. The fitted scale head is compared with original U-TV c=1, two all-category/CIGS-only validation constants and the per-datum posthoc oracle. Improvements in global PSNR must NOT be described as successful crack recovery unless line contrast and continuity also improve.}
\end{figure}
\clearpage

\section{Complete visual atlas of the twelve final-test reconstructions}

Each test image and each final noise realisation has two panels: the full 256-by-256 reconstruction and an enlarged region centred on the documented thin line. All methods within a figure use the SAME measurements and colour range. The source float32 reconstructions are preserved and the PDF panels are display-only.

\subsection{CIGS 0674}

Across the fresh test noise seeds: baseline PSNR 28.537 dB, global c PSNR 32.786 dB, Approach A PSNR 32.786 dB, posthoc oracle PSNR 32.786 dB. The learned head adds 0.000 dB over the global constant on this image. Line reconstruction quality must be inspected independently.

\begin{figure}[p]
\centering
\includegraphics[width=0.98\textwidth,height=0.78\textheight,keepaspectratio]{\detokenize{figures/CIGS_0674_noise5101_full.png}}
\caption{CIGS 0674, final noise seed 5101: whole image. Fixed fresh noise seed 5101: c=1 gives 28.526 dB; the learned scale 8.0000 gives 32.802 dB, and test oracle 8.0000 gives 32.802 dB. The oracle uses GT AFTER the model was selected and is not available at deployment.}
\end{figure}
\clearpage

\begin{figure}[p]
\centering
\includegraphics[width=0.98\textwidth,height=0.78\textheight,keepaspectratio]{\detokenize{figures/CIGS_0674_noise5101_roi.png}}
\caption{CIGS 0674, final noise seed 5101: thin-line crop. Fixed fresh noise seed 5101: c=1 gives 28.526 dB; the learned scale 8.0000 gives 32.802 dB, and test oracle 8.0000 gives 32.802 dB. The oracle uses GT AFTER the model was selected and is not available at deployment. Learned-scale signed retained line contrast is 0.015; a near-zero value indicates suppression, regardless of good global PSNR.}
\end{figure}
\clearpage

\begin{figure}[p]
\centering
\includegraphics[width=0.98\textwidth,height=0.78\textheight,keepaspectratio]{\detokenize{figures/CIGS_0674_noise5102_full.png}}
\caption{CIGS 0674, final noise seed 5102: whole image. Fixed fresh noise seed 5102: c=1 gives 28.492 dB; the learned scale 8.0000 gives 32.848 dB, and test oracle 8.0000 gives 32.848 dB. The oracle uses GT AFTER the model was selected and is not available at deployment.}
\end{figure}
\clearpage

\begin{figure}[p]
\centering
\includegraphics[width=0.98\textwidth,height=0.78\textheight,keepaspectratio]{\detokenize{figures/CIGS_0674_noise5102_roi.png}}
\caption{CIGS 0674, final noise seed 5102: thin-line crop. Fixed fresh noise seed 5102: c=1 gives 28.492 dB; the learned scale 8.0000 gives 32.848 dB, and test oracle 8.0000 gives 32.848 dB. The oracle uses GT AFTER the model was selected and is not available at deployment. Learned-scale signed retained line contrast is 0.056; a near-zero value indicates suppression, regardless of good global PSNR.}
\end{figure}
\clearpage

\begin{figure}[p]
\centering
\includegraphics[width=0.98\textwidth,height=0.78\textheight,keepaspectratio]{\detokenize{figures/CIGS_0674_noise5103_full.png}}
\caption{CIGS 0674, final noise seed 5103: whole image. Fixed fresh noise seed 5103: c=1 gives 28.594 dB; the learned scale 8.0000 gives 32.708 dB, and test oracle 8.0000 gives 32.708 dB. The oracle uses GT AFTER the model was selected and is not available at deployment.}
\end{figure}
\clearpage

\begin{figure}[p]
\centering
\includegraphics[width=0.98\textwidth,height=0.78\textheight,keepaspectratio]{\detokenize{figures/CIGS_0674_noise5103_roi.png}}
\caption{CIGS 0674, final noise seed 5103: thin-line crop. Fixed fresh noise seed 5103: c=1 gives 28.594 dB; the learned scale 8.0000 gives 32.708 dB, and test oracle 8.0000 gives 32.708 dB. The oracle uses GT AFTER the model was selected and is not available at deployment. Learned-scale signed retained line contrast is 0.059; a near-zero value indicates suppression, regardless of good global PSNR.}
\end{figure}
\clearpage

\subsection{CIGS 0985}

Across the fresh test noise seeds: baseline PSNR 28.047 dB, global c PSNR 31.863 dB, Approach A PSNR 31.863 dB, posthoc oracle PSNR 31.863 dB. The learned head adds 0.000 dB over the global constant on this image. Line reconstruction quality must be inspected independently.

\begin{figure}[p]
\centering
\includegraphics[width=0.98\textwidth,height=0.78\textheight,keepaspectratio]{\detokenize{figures/CIGS_0985_noise5101_full.png}}
\caption{CIGS 0985, final noise seed 5101: whole image. Fixed fresh noise seed 5101: c=1 gives 28.120 dB; the learned scale 8.0000 gives 31.890 dB, and test oracle 8.0000 gives 31.890 dB. The oracle uses GT AFTER the model was selected and is not available at deployment.}
\end{figure}
\clearpage

\begin{figure}[p]
\centering
\includegraphics[width=0.98\textwidth,height=0.78\textheight,keepaspectratio]{\detokenize{figures/CIGS_0985_noise5101_roi.png}}
\caption{CIGS 0985, final noise seed 5101: thin-line crop. Fixed fresh noise seed 5101: c=1 gives 28.120 dB; the learned scale 8.0000 gives 31.890 dB, and test oracle 8.0000 gives 31.890 dB. The oracle uses GT AFTER the model was selected and is not available at deployment. Learned-scale signed retained line contrast is 0.787; a near-zero value indicates suppression, regardless of good global PSNR.}
\end{figure}
\clearpage

\begin{figure}[p]
\centering
\includegraphics[width=0.98\textwidth,height=0.78\textheight,keepaspectratio]{\detokenize{figures/CIGS_0985_noise5102_full.png}}
\caption{CIGS 0985, final noise seed 5102: whole image. Fixed fresh noise seed 5102: c=1 gives 28.140 dB; the learned scale 8.0000 gives 31.945 dB, and test oracle 8.0000 gives 31.945 dB. The oracle uses GT AFTER the model was selected and is not available at deployment.}
\end{figure}
\clearpage

\begin{figure}[p]
\centering
\includegraphics[width=0.98\textwidth,height=0.78\textheight,keepaspectratio]{\detokenize{figures/CIGS_0985_noise5102_roi.png}}
\caption{CIGS 0985, final noise seed 5102: thin-line crop. Fixed fresh noise seed 5102: c=1 gives 28.140 dB; the learned scale 8.0000 gives 31.945 dB, and test oracle 8.0000 gives 31.945 dB. The oracle uses GT AFTER the model was selected and is not available at deployment. Learned-scale signed retained line contrast is 0.830; a near-zero value indicates suppression, regardless of good global PSNR.}
\end{figure}
\clearpage

\begin{figure}[p]
\centering
\includegraphics[width=0.98\textwidth,height=0.78\textheight,keepaspectratio]{\detokenize{figures/CIGS_0985_noise5103_full.png}}
\caption{CIGS 0985, final noise seed 5103: whole image. Fixed fresh noise seed 5103: c=1 gives 27.881 dB; the learned scale 8.0000 gives 31.755 dB, and test oracle 8.0000 gives 31.755 dB. The oracle uses GT AFTER the model was selected and is not available at deployment.}
\end{figure}
\clearpage

\begin{figure}[p]
\centering
\includegraphics[width=0.98\textwidth,height=0.78\textheight,keepaspectratio]{\detokenize{figures/CIGS_0985_noise5103_roi.png}}
\caption{CIGS 0985, final noise seed 5103: thin-line crop. Fixed fresh noise seed 5103: c=1 gives 27.881 dB; the learned scale 8.0000 gives 31.755 dB, and test oracle 8.0000 gives 31.755 dB. The oracle uses GT AFTER the model was selected and is not available at deployment. Learned-scale signed retained line contrast is 0.887; a near-zero value indicates suppression, regardless of good global PSNR.}
\end{figure}
\clearpage

\subsection{CIGS 0917}

Across the fresh test noise seeds: baseline PSNR 28.078 dB, global c PSNR 32.115 dB, Approach A PSNR 32.115 dB, posthoc oracle PSNR 32.115 dB. The learned head adds 0.000 dB over the global constant on this image. Line reconstruction quality must be inspected independently.

\begin{figure}[p]
\centering
\includegraphics[width=0.98\textwidth,height=0.78\textheight,keepaspectratio]{\detokenize{figures/CIGS_0917_noise5101_full.png}}
\caption{CIGS 0917, final noise seed 5101: whole image. Fixed fresh noise seed 5101: c=1 gives 28.180 dB; the learned scale 8.0000 gives 32.193 dB, and test oracle 8.0000 gives 32.193 dB. The oracle uses GT AFTER the model was selected and is not available at deployment.}
\end{figure}
\clearpage

\begin{figure}[p]
\centering
\includegraphics[width=0.98\textwidth,height=0.78\textheight,keepaspectratio]{\detokenize{figures/CIGS_0917_noise5101_roi.png}}
\caption{CIGS 0917, final noise seed 5101: thin-line crop. Fixed fresh noise seed 5101: c=1 gives 28.180 dB; the learned scale 8.0000 gives 32.193 dB, and test oracle 8.0000 gives 32.193 dB. The oracle uses GT AFTER the model was selected and is not available at deployment. Learned-scale signed retained line contrast is 0.845; a near-zero value indicates suppression, regardless of good global PSNR.}
\end{figure}
\clearpage

\begin{figure}[p]
\centering
\includegraphics[width=0.98\textwidth,height=0.78\textheight,keepaspectratio]{\detokenize{figures/CIGS_0917_noise5102_full.png}}
\caption{CIGS 0917, final noise seed 5102: whole image. Fixed fresh noise seed 5102: c=1 gives 28.020 dB; the learned scale 8.0000 gives 32.201 dB, and test oracle 8.0000 gives 32.201 dB. The oracle uses GT AFTER the model was selected and is not available at deployment.}
\end{figure}
\clearpage

\begin{figure}[p]
\centering
\includegraphics[width=0.98\textwidth,height=0.78\textheight,keepaspectratio]{\detokenize{figures/CIGS_0917_noise5102_roi.png}}
\caption{CIGS 0917, final noise seed 5102: thin-line crop. Fixed fresh noise seed 5102: c=1 gives 28.020 dB; the learned scale 8.0000 gives 32.201 dB, and test oracle 8.0000 gives 32.201 dB. The oracle uses GT AFTER the model was selected and is not available at deployment. Learned-scale signed retained line contrast is 0.872; a near-zero value indicates suppression, regardless of good global PSNR.}
\end{figure}
\clearpage

\begin{figure}[p]
\centering
\includegraphics[width=0.98\textwidth,height=0.78\textheight,keepaspectratio]{\detokenize{figures/CIGS_0917_noise5103_full.png}}
\caption{CIGS 0917, final noise seed 5103: whole image. Fixed fresh noise seed 5103: c=1 gives 28.034 dB; the learned scale 8.0000 gives 31.951 dB, and test oracle 8.0000 gives 31.951 dB. The oracle uses GT AFTER the model was selected and is not available at deployment.}
\end{figure}
\clearpage

\begin{figure}[p]
\centering
\includegraphics[width=0.98\textwidth,height=0.78\textheight,keepaspectratio]{\detokenize{figures/CIGS_0917_noise5103_roi.png}}
\caption{CIGS 0917, final noise seed 5103: thin-line crop. Fixed fresh noise seed 5103: c=1 gives 28.034 dB; the learned scale 8.0000 gives 31.951 dB, and test oracle 8.0000 gives 31.951 dB. The oracle uses GT AFTER the model was selected and is not available at deployment. Learned-scale signed retained line contrast is 0.804; a near-zero value indicates suppression, regardless of good global PSNR.}
\end{figure}
\clearpage

\subsection{CIGS 0864}

Across the fresh test noise seeds: baseline PSNR 28.330 dB, global c PSNR 32.152 dB, Approach A PSNR 32.082 dB, posthoc oracle PSNR 32.152 dB. The learned head adds -0.070 dB over the global constant on this image. Line reconstruction quality must be inspected independently.

\begin{figure}[p]
\centering
\includegraphics[width=0.98\textwidth,height=0.78\textheight,keepaspectratio]{\detokenize{figures/CIGS_0864_noise5101_full.png}}
\caption{CIGS 0864, final noise seed 5101: whole image. Fixed fresh noise seed 5101: c=1 gives 28.336 dB; the learned scale 8.0000 gives 32.147 dB, and test oracle 8.0000 gives 32.147 dB. The oracle uses GT AFTER the model was selected and is not available at deployment.}
\end{figure}
\clearpage

\begin{figure}[p]
\centering
\includegraphics[width=0.98\textwidth,height=0.78\textheight,keepaspectratio]{\detokenize{figures/CIGS_0864_noise5101_roi.png}}
\caption{CIGS 0864, final noise seed 5101: thin-line crop. Fixed fresh noise seed 5101: c=1 gives 28.336 dB; the learned scale 8.0000 gives 32.147 dB, and test oracle 8.0000 gives 32.147 dB. The oracle uses GT AFTER the model was selected and is not available at deployment. Learned-scale signed retained line contrast is 0.840; a near-zero value indicates suppression, regardless of good global PSNR.}
\end{figure}
\clearpage

\begin{figure}[p]
\centering
\includegraphics[width=0.98\textwidth,height=0.78\textheight,keepaspectratio]{\detokenize{figures/CIGS_0864_noise5102_full.png}}
\caption{CIGS 0864, final noise seed 5102: whole image. Fixed fresh noise seed 5102: c=1 gives 28.248 dB; the learned scale 11.3137 gives 31.905 dB, and test oracle 8.0000 gives 32.116 dB. The oracle uses GT AFTER the model was selected and is not available at deployment.}
\end{figure}
\clearpage

\begin{figure}[p]
\centering
\includegraphics[width=0.98\textwidth,height=0.78\textheight,keepaspectratio]{\detokenize{figures/CIGS_0864_noise5102_roi.png}}
\caption{CIGS 0864, final noise seed 5102: thin-line crop. Fixed fresh noise seed 5102: c=1 gives 28.248 dB; the learned scale 11.3137 gives 31.905 dB, and test oracle 8.0000 gives 32.116 dB. The oracle uses GT AFTER the model was selected and is not available at deployment. Learned-scale signed retained line contrast is 0.861; a near-zero value indicates suppression, regardless of good global PSNR.}
\end{figure}
\clearpage

\begin{figure}[p]
\centering
\includegraphics[width=0.98\textwidth,height=0.78\textheight,keepaspectratio]{\detokenize{figures/CIGS_0864_noise5103_full.png}}
\caption{CIGS 0864, final noise seed 5103: whole image. Fixed fresh noise seed 5103: c=1 gives 28.407 dB; the learned scale 8.0000 gives 32.194 dB, and test oracle 8.0000 gives 32.194 dB. The oracle uses GT AFTER the model was selected and is not available at deployment.}
\end{figure}
\clearpage

\begin{figure}[p]
\centering
\includegraphics[width=0.98\textwidth,height=0.78\textheight,keepaspectratio]{\detokenize{figures/CIGS_0864_noise5103_roi.png}}
\caption{CIGS 0864, final noise seed 5103: thin-line crop. Fixed fresh noise seed 5103: c=1 gives 28.407 dB; the learned scale 8.0000 gives 32.194 dB, and test oracle 8.0000 gives 32.194 dB. The oracle uses GT AFTER the model was selected and is not available at deployment. Learned-scale signed retained line contrast is 0.951; a near-zero value indicates suppression, regardless of good global PSNR.}
\end{figure}
\clearpage

\section{Limitations and scientific interpretation}

\begin{enumerate}

\item The 4 images were already examined during exploratory Stage I, so Stage II is not fully blind

\item CIGS 0674 may be from the old map-network training partition; excluded from scale head but pipeline-level unseen-image claim is invalid

\item A 50-image dataset is a pilot, not an established generalisation benchmark

\item Global quality-optimal scaling may destroy material cracks

\item Oracle scale labels are computed on a finite 19-point grid. Their best value depends on GT, and per-datum predictability is not guaranteed. A wide near-optimal plateau can make hard-label regression unnecessarily unstable.

\item The scale correction multiplies the full map and cannot alter its spatial shape. Stronger global TV may improve PSNR while deleting a one-pixel defect.

\end{enumerate}

\section{Reproducibility and numerical integrity}

The scripts refuse incomplete Stage I/II data, fail if source images change, freeze image partition and pretrained U-TV checkpoint hashes, and select the scale predictor on validation data before evaluating the new test noise seeds. The 50 training images come from original training IDs 0000--0799 per category; the 25 validation images from 0800--0899 per category. All four CIGS test IDs are excluded from both.

\begin{center}
\small
\begin{tabular}{lr}
\toprule
Artifact & Recorded value\\
\midrule
Experiment & Experiment 1 on 28-10-2026\\
Map model checkpoint & best\_val\_loss\_s0.20\_cJ6.pt\\
Map checkpoint SHA256 & 4da3f177ece6dc6f23979a1d...\\
Scale-head SHA256 & 780c658259aee1d1d5b4b7a3...\\
LION Git revision & ce5dcb4780780d9779eee6c4...\\
Training/validation sample manifest SHA256 & 8732a038051159a74d2ef8a2...\\
Stage I run root & /store/LION/ttv22/pcm/results/experiment\_1\_on\_28-10-2026/old\\
Stage II run root & /store/LION/ttv22/pcm/results/experiment\_1\_on\_28-10-2026/approach\_a/old\\
\bottomrule
\end{tabular}
\end{center}

\section{Conclusions and suggested follow-up}

The central experimental question is whether a per-datum scale predictor captures a meaningful additional improvement beyond the BEST single validation-calibrated multiplier. On these four CIGS images it achieved -0.018 dB of extra mean PSNR compared with the global scale; the oracle gap of 0.000 dB is the posthoc headroom, not a guaranteed achievable gain. If the learned gain is small, prefer a single calibrated constant. If the global metric gain is large but thin lines disappear, scalar calibration alone is not sufficient for defect imaging without independently validated crack masks; investigate line-aware objectives or directional regularisation. A future genuinely blinded test should use NEW image identities never used to design the experiment.

\end{document}
